%% LyX 2.4.4 created this file.  For more info, see https://www.lyx.org/.
%% Do not edit unless you really know what you are doing.
\documentclass[english]{article}
\usepackage[T1]{fontenc}
\usepackage[latin9]{inputenc}
\usepackage{float}
\usepackage{units}
\usepackage{enumitem}
\usepackage{amsmath}
\usepackage{amsthm}
\usepackage{graphicx}
\usepackage{geometry}
\geometry{verbose,tmargin=1in,bmargin=1in,lmargin=1in,rmargin=1in}

\makeatletter
%%%%%%%%%%%%%%%%%%%%%%%%%%%%%% Textclass specific LaTeX commands.
\numberwithin{equation}{section}
\numberwithin{figure}{section}
\newlength{\lyxlabelwidth}      % auxiliary length 

%%%%%%%%%%%%%%%%%%%%%%%%%%%%%% User specified LaTeX commands.
\usepackage{fancyhdr}
\usepackage{lastpage}
\pagestyle{fancy}
\usepackage{enumitem}
\usepackage{ifthen}
%sets math fonts to be more readable
\everymath=\expandafter{\the\everymath\displaystyle}
%\DeclareMathSizes{10}{10}{10}{10}
\usepackage{multicol}
% Added by lyx2lyx
\setlength{\parskip}{\smallskipamount}
\setlength{\parindent}{0pt}

\makeatother

\usepackage{babel}
\begin{document}
\renewcommand{\author}{Dr.\,James Ashe}

\newcommand{\classNum}{331}

\newcommand{\classSec}{A}

\newcommand{\nameType}{Name:}

\newcommand{\examType}{Exam 2 Make-up}

\renewcommand{\date}{November 19, 2012}

%%First page's header%%%%%%%%%%%%%%%%%%%%%%
\fancypagestyle{firststyle} %%header settings for first pagr
{
	\fancyhead[L]{MTH \classNum{}(\classSec{}) \author{}\\
	\textbf{\examType{}} \date{}}
	\fancyhead[C]{\textbf{\nameType}}
	\setlength{\headsep}{14pt}
}
%%%%%%%%%%%%%%%%%%%%%%%%%%%%%%%%%%%%%%%%%%

%%Next page's header%%%%%%%%%%%%%%%%%%%%%%%
\setlist{nolistsep}
\lhead{\classNum{}(\classSec{})} 
\chead{\textbf{\nameType}} 
\cfoot{\thepage{}~of~\pageref{LastPage}} 
\setlength{\headsep}{5pt} 
%%%%%%%%%%%%%%%%%%%%%%%%%%%%%%%%%%%%%%%%%%%

%%Various settings; see the preamble for other settings%%%%%
%%\setenumerate[0]{leftmargin=12pt,itemindent=0pt,labelwidth=0pt}
\renewcommand{\labelenumi}{\textbf{\arabic{enumi}.}}
\renewcommand{\labelenumii}{\textbf{\alph{enumii}.}}
\thispagestyle{firststyle}
%%%%%%%%%%%%%%%%%%%%%%%%%%%%%%%%%%%%%%%%%%
\begin{enumerate}[resume]
\item Let $f(x)=2e^{x}$.

\begin{enumerate}[resume]
\item Find the Taylor polynomial of order $n=3$ for $f(x)$ centered at
$0$.
\item Use the Taylor polynomial from above to approximate $2e^{-0.06}$.
\item Compute the absolute error in the approximation assuming the exact
value is given by a calculator.\vspace{3cm}\vfill
\end{enumerate}
\end{enumerate}
\emph{Determine the radius of convergence of the following power series. }
\begin{enumerate}[resume]
\item ${\displaystyle \sum\frac{x^{k}}{0.7^{k}}}$\vfill
\end{enumerate}
\emph{Use the geometric series ${\displaystyle \sum_{k=0}^{\infty}x^{k}=\frac{1}{1-x}}$,
for $\left|x\right|<1$ to find a power series representations for
the following functions, and determine the interval of convergence.}
\begin{enumerate}[resume]
\item $f(x)=\frac{5}{1+x}$.\vfill
\item $f(x)=\frac{2x}{1-x^{2}}$.\vfill\newpage
\end{enumerate}
\emph{Eliminate the parameter in following equations to obtain a single
equation in terms of $x$ and $y$. }
\begin{enumerate}[resume]
\item $x=(t+1)^{2}$, $y=t+2$.\vfill
\end{enumerate}
\emph{Find an equation of the line tangent to the following curve
at the given point $t$. }
\begin{enumerate}[resume]
\item $x=3\sin t$, $y=3\cos t$; $t=\nicefrac{\pi}{2}$.\vfill
\end{enumerate}
\emph{Give a set of parametric equations that describes the curve.}
\begin{enumerate}[resume]
\item A bicyclist rides clockwise with a constant speed around a cricular
track with radius 50 meters, completing one lap in 30 seconds.\vfill\newpage
\end{enumerate}
\emph{Locate the following polar coordinates on the graph. }
\begin{enumerate}[resume]
\item A:$\left(1.5,\frac{\pi}{6}\right)$, B:$\left(-1.5,\frac{\pi}{6}\right)$,
C:$\left(1.5,\frac{\pi}{6}+3\pi\right)$, D:$\left(1.5,\frac{\pi}{6}-8\pi\right)$
\begin{figure}[H]
\begin{centering}
\includegraphics[width=2in]{D:/home/jim/JamesAshe/JCSU/Calc_3/slides_Calc3/images/polar_graph_-2-2a.png}
\par\end{centering}
\caption{Use this graph to plot the polar coordinates.}
\end{figure}
\end{enumerate}
\emph{Express the following }\textbf{\emph{polar}}\emph{ coordinates
in }\textbf{\emph{Cartesian}}\emph{ coordinates.}
\begin{enumerate}[resume]
\item $\left(4,5\pi\right)$\vfill
\end{enumerate}
\emph{Express the following }\textbf{\emph{Cartesian}}\emph{ coordinates
in }\textbf{\emph{polar}}\emph{ coordinates in at least }\textbf{\emph{two}}\emph{
different ways. }
\begin{enumerate}[resume]
\item $\left(-1,\sqrt{3}\right)$\vfill\newpage
\end{enumerate}
\emph{Multiple Choice: Identify which set of parametric curves describe
the graph.}
\begin{enumerate}[resume]
\item Parametric equations that describe the curve are

\begin{enumerate}[resume]
\item [(A)]$x=\cos t$, $y=\sin t$, $0\leq t\leq2\pi$.
\item [(B)]$x=\cos t$, $y=\sin2t$, $-\pi\leq t\leq\pi$.
\item [(C)]$x=\cos2t$, $y=\sin t$, $-\pi\leq t\leq\pi$.
\begin{figure}[H]
\centering{}\includegraphics[width=3cm]{slides_Calc3/images/cos2t_sint.png}
\end{figure}
\end{enumerate}
\end{enumerate}
\emph{Find the slope of the line tangent to the following polar curves
at the given point using calculus. }
\begin{enumerate}[resume]
\item $r=2\cos2\theta$; $\left(2,\frac{\pi}{2}\right)$.\vfill
\end{enumerate}
\emph{Find the area of the region using calculus. }
\begin{enumerate}[resume]
\item The region enclosed by $r=2+\cos4\theta$.\vfill
\end{enumerate}

\end{document}
